%%% تنظیمات مربوط به بسته  glossaries
%%% تعریف استایل برای واژه‌نامه فارسی به انگلیسی، در این استایل واژه‌های فارسی در سمت راست و واژه‌های انگلیسی در سمت چپ خواهند آمد. از حالت گروه ‌بندی استفاده می‌کنیم، 
%%% یعنی واژه‌ها در گروه‌هایی به ترتیب حروف الفبا مرتب می‌شوند، مثلا:
%%% الف
%%% افتصاد ................................... Economy
%%% اشکال ........................................ Failure
%%% ش
%%% شبکه ...................................... Network
\newglossarystyle{myFaToEn}{%
	\renewenvironment{theglossary}{}{}
	\renewcommand*{\glsgroupskip}{\vskip 10mm}
	\renewcommand*{\glsgroupheading}[1]{\subsection*{\glsgetgrouptitle{##1}}}
	\renewcommand*{\glossentry}[2]{\noindent\glsentryname{##1}\dotfill\space \glsentrytext{##1}
		
	}
}

%% % تعریف استایل برای واژه‌نامه انگلیسی به فارسی، در این استایل واژه‌های فارسی در سمت راست و واژه‌های انگلیسی در سمت چپ خواهند آمد. از حالت گروه ‌بندی استفاده می‌کنیم، 
%% % یعنی واژه‌ها در گروه‌هایی به ترتیب حروف الفبا مرتب می‌شوند، مثلا:
%% % E
%%% Economy ............................... اقتصاد
%% % F
%% % Failure................................... اشکال
%% %N
%% % Network ................................. شبکه

\newglossarystyle{myEntoFa}{%
	%%% این دستور در حقیقت عملیات گروه‌بندی را انجام می‌دهد. بدین صورت که واژه‌ها در بخش‌های جداگانه گروه‌بندی می‌شوند، 
	%%% عنوان بخش همان نام حرفی است که هر واژه در آن گروه با آن شروع شده است. 
	\renewenvironment{theglossary}{}{}
	\renewcommand*{\glsgroupskip}{\vskip 10mm}
	\renewcommand*{\glsgroupheading}[1]{\begin{LTR} \subsection*{\glsgetgrouptitle{##1}} \end{LTR}}
	%%% در این دستور نحوه نمایش واژه‌ها می‌آید. در این جا واژه فارسی در سمت راست و واژه انگلیسی در سمت چپ قرار داده شده است، و بین آن با نقطه پر می‌شود. 
	\renewcommand*{\glossentry}[2]{\noindent\glsentrytext{##1}\dotfill\space \glsentryname{##1}
		
	}
}

%%% تعیین استایل برای فهرست اختصارات
\newglossarystyle{myAbbrlist}{%
	%%% این دستور در حقیقت عملیات گروه‌بندی را انجام می‌دهد. بدین صورت که اختصارات‌ در بخش‌های جداگانه گروه‌بندی می‌شوند، 
	%%% عنوان بخش همان نام حرفی است که هر اختصار در آن گروه با آن شروع شده است. 
	\renewenvironment{theglossary}{}{}
	\renewcommand*{\glsgroupskip}{\vskip 10mm}
	\renewcommand*{\glsgroupheading}[1]{\begin{LTR} \subsection*{\glsgetgrouptitle{##1}} \end{LTR}}
	%%% در این دستور نحوه نمایش اختصارات می‌آید. در این جا حالت کوچک اختصار در سمت چپ و حالت بزرگ در سمت راست قرار داده شده است، و بین آن با نقطه پر می‌شود. 
	\renewcommand*{\glossentry}[2]{\noindent\Glsentrylong{##1}\dotfill\space \glsentrytext{##1} 
		
	}
	%%% تغییر نام محیط abbreviation به فهرست اختصارات
	\renewcommand*{\acronymname}{\rl{فهرست اختصارات}}
}

%%% برای اجرا xindy بر روی فایل .tex و تولید واژه‌نامه‌ها و فهرست اختصارات و فهرست نمادها یکسری  فایل تعریف شده است.‌ Latex داده های مربوط به واژه‌نامه و .. را در این 
%%%  فایل‌ها نگهداری می‌کند. مهم‌ترین option‌ این قسمت این است که 
%%% عنوان واژه‌نامه‌ها و یا فهرست اختصارات و یا فهرست نمادها را می‌توانید در این‌جا مشخص کنید. 
%%% در این جا عباراتی مثل glg، gls، glo و ... پسوند فایل‌هایی است که برای xindy بکار می‌روند. 
\newglossary[glg]{english}{gls}{glo}{واژه‌نامهٔ انگلیسی به فارسی}
\newglossary[blg]{persian}{bls}{blo}{واژه‌نامهٔ فارسی به انگلیسی}
\makeglossaries
\glsdisablehyper
%%% تعاریف مربوط به تولید واژه‌نامه و فهرست اختصارات و فهرست نمادها
%%%  در این فایل یکسری دستورات عمومی برای وارد کردن واژه‌نامه آمده است.
%%%  به دلیل این‌که قرار است این دستورات پایه‌ای را بازنویسی کنیم در این‌جا تعریف می‌کنیم. 
\let\oldgls\gls
\let\oldglspl\glspl

\makeatletter

\renewrobustcmd*{\gls}{\@ifstar\@msgls\@mgls}
\newcommand*{\@mgls}[1] {\ifthenelse{\equal{\glsentrytype{#1}}{english}}{\oldgls{#1}\glsuseri{f-#1}}{\lr{\oldgls{#1}}}}
\newcommand*{\@msgls}[1]{\ifthenelse{\equal{\glsentrytype{#1}}{english}}{\glstext{#1}\glsuseri{f-#1}}{\lr{\glsentryname{#1}}}}

\renewrobustcmd*{\glspl}{\@ifstar\@msglspl\@mglspl}
\newcommand*{\@mglspl}[1] {\ifthenelse{\equal{\glsentrytype{#1}}{english}}{\oldglspl{#1}\glsuseri{f-#1}}{\oldglspl{#1}}}
\newcommand*{\@msglspl}[1]{\ifthenelse{\equal{\glsentrytype{#1}}{english}}{\glsplural{#1}\glsuseri{f-#1}}{\glsentryplural{#1}}}

\makeatother

\newcommand{\newword}[4]{
	\newglossaryentry{#1}     {type={english},name={\lr{#2}},plural={#4},text={#3},description={}}
	\newglossaryentry{f-#1} {type={persian},name={#3},text={\lr{#2}},description={}}
}

%%% بر طبق این دستور، در اولین باری که واژه مورد نظر از واژه‌نامه وارد شود، پاورقی زده می‌شود. 
\defglsentryfmt[english]{\glsgenentryfmt\ifglsused{\glslabel}{}{\LTRfootnote{\glsentryname{\glslabel}}}}

%%% بر طبق این دستور، در اولین باری که واژه مورد نظر از فهرست اختصارات وارد شود، پاورقی زده می‌شود. 
\defglsentryfmt[acronym]{\glsentryname{\glslabel}\ifglsused{\glslabel}{}{\footnote{\glsentrydesc{\glslabel}}}}


%%%%%% ============================================================================================================

%%============================ دستور برای قرار دادن فهرست اختصارات 
\newcommand{\printabbreviation}{
	%\cleardoublepage
	%\phantomsection
	\baselineskip=.75cm
	\setglossarystyle{myAbbrlist}
	%\begin{LTR}
	\Oldprintglossary[type=acronym]	
	%\end{LTR}
	\clearpage
}%

\newcommand{\printacronyms}{\printabbreviation}

% فهرست اختصارات فنی پایان‌نامه.  در این پروژه اختصارات در متن با پاورقی
% انگلیسی معرفی می‌شوند؛ جدول زیر شکل ثابت و قابل‌خواندن فهرست را مستقل از
% مرتب‌سازی xindy نگه می‌دارد.
\newcommand{\printabbreviationmanual}{%
\clearpage
\chapter*{\rl{فهرست اختصارات}}
\addcontentsline{toc}{chapter}{\rl{فهرست اختصارات}}
\small
\begin{longtable}{@{}>{\raggedleft\arraybackslash}p{0.31\textwidth} >{\centering\arraybackslash}p{0.13\textwidth} >{\raggedright\arraybackslash}p{0.46\textwidth}@{}}
\toprule
معادل فارسی & اختصار & صورت کامل انگلیسی \\
\midrule
\endhead
بینایی ماشین & \lr{CV} & \lr{Computer Vision} \\
تخمین عمق تک‌دوربینه & \lr{MDE} & \lr{Monocular Depth Estimation} \\
شبکهٔ عصبی کانولوشنی & \lr{CNN} & \lr{Convolutional Neural Network} \\
میدان دید & \lr{FoV} & \lr{Field of View} \\
سیستم‌های کمک‌رانندهٔ پیشرفته & \lr{ADAS} & \lr{Advanced Driver-Assistance Systems} \\
ترنسفورمر بینایی & \lr{ViT} & \lr{Vision Transformer} \\
ترنسفورمر پیش‌بینی چگال & \lr{DPT} & \lr{Dense Prediction Transformer} \\
قرمز--سبز--آبی & \lr{RGB} & \lr{Red--Green--Blue} \\
قرمز--سبز--آبی و عمق & \lr{RGB-D} & \lr{Red--Green--Blue--Depth} \\
پرسپترون چندلایه & \lr{MLP} & \lr{Multi-Layer Perceptron} \\
تعدیل خطی در سطح ویژگی & \lr{FiLM} & \lr{Feature-wise Linear Modulation} \\
واحد خطی خطای گاوسی & \lr{GELU} & \lr{Gaussian Error Linear Unit} \\
واحد خطی یکسوشده & \lr{ReLU} & \lr{Rectified Linear Unit} \\
واحد پردازش گرافیکی & \lr{GPU} & \lr{Graphics Processing Unit} \\
قالب تبادل اطلاعات تصویری & \lr{EXIF} & \lr{Exchangeable Image File Format} \\
ریشهٔ میانگین مربعات خطا & \lr{RMSE} & \lr{Root Mean Squared Error} \\
خطای نسبی مطلق & \lr{AbsRel} & \lr{Absolute Relative Error} \\
خطای نسبی مربعی & \lr{SqRel} & \lr{Squared Relative Error} \\
تابع هزینهٔ لگاریتمی تغییرناپذیر نسبت به مقیاس & \lr{SiLog} & \lr{Scale-Invariant Logarithmic Loss} \\
عمق متریک صفرنمونه & \lr{DMD} & \lr{Zero-Shot Metric Depth} \\
\bottomrule
\end{longtable}
\normalsize
\clearpage
}%
%%% در این جا محیط هر دو واژه‌نامه را باز تعریف کرده ایم، تا اولا مشکل قرار دادن صفحه اضافی را حل کنیم، ثانیا عنوان واژه‌نامه ها را با دستور addcontentlist وارد فهرست مطالب کرده ایم.
\let\Oldprintglossary\printglossary
\renewcommand{\printglossary}{
	\let\appendix\relax
	%% تنظیم کننده فاصله بین خطوط در این قسمت
	\clearpage
	%\phantomsection
	%% این دستور موجب این می‌شود که واژه‌نامه‌ها در  حالت دو ستونی نوشته شود. 
	\twocolumn{}
	\setglossarystyle{myFaToEn}
	\Oldprintglossary[type=persian]
	\clearpage
	%\phantomsection
	\setglossarystyle{myEntoFa}
	\Oldprintglossary[type=english]	
	\onecolumn{}
}%
%%%%%%
